\documentclass[../libro.tex]{subfiles}

\begin{document}

\ifSubfilesClassLoaded{\frontmatter}{}

\chapter{序}

这是一本关于行列式的书.
行列式是一个初等的, 有用的工具.
行列式在现在不如在过去那么重要了,
但学习行列式不是不好的.
行列式在代数, 微积分, 几何是有用的.
于是, 我写了本书.

我假定您有中学代数基础
(包括但不限于可比数 (\angla{rational numbers}),
实数,
代数式,
\(2\)~元~\(1\)~次方程组,
整式的加减乘,
因式分解,
分式,
\(1\)~元~\(1\)~次不等式,
常量与变量,
函数,
\(2\)~次根式,
\(1\)~元~\(2\)~次方程;
我列不完).
我希望本书不会是难的,
% 但因为我几乎无法作实验
% (用我的书教初学者行列式),
但因为我几乎无法作实验,
故我只能希望如此.

本书有一些创新, 但不多, 且不重要.

本书的最新版本在
\url{https://www.bananaspace.org/wiki/讲义:行列式入门}
与
\url{https://github.com/septsea/det}.

若您发现本书的错误, 请使我知道它.
谢谢.

\vspace{4ex}\inscribe{\laauxtoro}{\ladato}

\end{document}

Well, the book was intended
to be a book for beginners,
but later I began to suspect
that the book was not really suitable for beginners.
So I rewrote the preface.
Here is the old preface:

\begin{comment}
本书的标题是《\latitolo{}》
(\textit{Enkonduko al determinantoj}).
您按字面意思理解此标题就好:
我
(即作者, 下同)
带您入 ``行列式之门''
(la pordo al determinantoj).

本书是一本为初学者准备的行列式教材.
行列式是一个有用的工具.
我认为, 学习此工具是有用的.
本书用简单的归纳法定义行列式,
并证明了关于行列式的一些结论.

本书的目标不是只教您行列式.
或许, 您知道,
中学数学里的几何学可练习我们的逻辑思维.
我希望本书也能有类似的作用.
具体地, 我希望本书也可练习您的逻辑思维.
我想, 逻辑思维对我们是好的,
即使我们不从事数学的工作.

我假定您学过中学数学;
或者, 我不假定您学过大学数学.
这其实对我是一个大挑战:
几年前, 我就不是一个中学生了.
并且, 我也没有学习教育的理论.
所以, 我不知道,
本书是否能使一个只有中学数学基础的人学明白行列式.
但我想挑战此事.
于是,
本书 (的初版) 就出生了.

本书的最新版本在
\url{https://www.bananaspace.org/wiki/讲义:行列式入门}
与
\url{https://github.com/septsea/det}.

我的文学不好.
那么, 我就说这么多.
\end{comment}
